% draft_co2_skeleton.tex  (v2, 2026-08-22 late)
% CO2-ONLY paper: scale-composition-technique framework, Feyrer IV primary.
% Supersedes v1 (trade-linked "carbon price" framing dropped per Sunbin).
% Numbers from: co2_feyrer_full.log, co2_main.log (tourism lens),
%   co2_airport_panel.log, co2_feyrer_sizecheck.log, co2_country_sizecheck2.log,
%   attribution one-off (356 Mt), 04/09 SAF arithmetic.
% Figures: CO2_map_levels_2023.png, CO2_map_mismatch_2023.png,
%   CO2_efficiency_curve.png, CO2_saf_scenarios.png (needs Mt-version rebuild),
%   CO2_rf_quintile.png (now appendix: flat under Feyrer).
% TODO blocks marked [TODO]; pending-coauthor [YIFU]/[FANGYU].
\documentclass[12pt]{article}
\usepackage{amsmath,graphicx,booktabs,threeparttable,float}
\newcommand{\sym}[1]{\rlap{$^{#1}$}}
\title{Scale Beats Efficiency: Global Air Connectivity and the Carbon
Footprint of Aviation}
\author{[author order TBD]}
\date{Skeleton v2, 22 August 2026}
\begin{document}
\maketitle

%=======================================================================
\section{Introduction}
%=======================================================================
% Skeleton chain:
% 1. Hook: aviation is the hardest-to-abate transport sector and the most
%    network-dependent one. Countries invest in connectivity; what does that
%    buy in carbon?
% 2. Theory gap: the emissions consequence of connectivity growth is
%    ambiguous ex ante. Growth adds traffic (scale), but hubbing consolidates
%    flying into larger aircraft and longer stages (technique), and shifts
%    the mix toward international long-haul (composition). Which margin wins
%    is an empirical question, and no causal estimate exists at global scale.
% 3. This paper: flight-stage-resolved CO2 for 6,000+ airports x 28 years
%    (EEA/EMEP methodology) x the GACI network panel, identified with a
%    Feyrer-type air/sea market-access instrument (KP F = 153).
% 4. Findings preview:
%    (i) scale wins: the emissions elasticity of connectivity is 5.7,
%        far above one (Wald p<0.001), identical across allocation rules;
%    (ii) technique offsets at the node: traffic rises faster than emissions
%        (6.1 vs 5.7), so CO2 per seat-km falls (-0.40 country, -0.49
%        airport, 284 -> 77 kg per 1,000 seat-km across the hub ladder);
%    (iii) attribution: connectivity growth since 1996 accounts for 42.5%
%        (356 Mt) of 2023 aviation CO2; China +88 Mt, USA +36, UAE +25,
%        Germany -16;
%    (iv) only fuel switching bends the curve: ReFuelEU-style SAF paths cut
%        the 2023 total from 838 Mt to 457-545 Mt by 2050 depending on
%        life-cycle savings.
% 5. Contributions: first causal connectivity-CO2 elasticity; the
%    scale-composition-technique (SCT) decomposition applied to aviation;
%    the attribution-versus-physical-emissions mismatch; SAF policy bridge.

%=======================================================================
\section{Conceptual framework and hypotheses}
%=======================================================================
% SCT decomposition (Grossman-Krueger 1995; Copeland-Taylor 2004) applied to
% aviation: E = seat-km (scale) x mix (composition) x CO2/seat-km (technique).
% Connectivity growth moves all three margins; the net elasticity is the
% empirical object. Outcome mapping: scale = ln seat-km; composition =
% international share / intl vs total emissions; technique = ln CO2/seat-km.
\paragraph{H1 (scale and proportionality).} Connectivity growth raises
aviation CO2; the elasticity exceeds one if scale dominates technique.
% Verdict: supported, beta = 5.67 (SE 0.41), Wald p(beta=1) < 0.001,
% robust across LTO/fuel-uplift/50-50 allocations (5.6-5.9).
\paragraph{H2 (technique).} Connectivity lowers emissions per seat-km through
hub consolidation, larger aircraft, and longer stage lengths.
% Verdict: supported at both levels: country -0.399 (SE 0.118), airport
% -0.494 (SE 0.044, within-country), efficiency curve 284 -> 77 kg.
\paragraph{H3 (composition).} Hub-quality gains internationalise the traffic
mix.
% Verdict: supported through the tourism-IV lens (total n.s. vs intl 3.16,
% SE 0.77) and the airport panel (intl share +0.33, SE 0.02).
\paragraph{H4 (attribution geography).} The emissions attributable to
connectivity growth are geographically concentrated and diverge from where
emissions physically occur, with accounting consequences for CORSIA and
bunker conventions.
% Verdict: supported: attributed growth concentrated in China/Gulf/Asia,
% declines in mature European hubs; attributed-minus-physical LTO share
% +1.4 pp (USA, UAE) vs -3.5 pp (China).

%=======================================================================
\section{Data}
%=======================================================================
\paragraph{Emissions.} Flight-level CO2 from OAG schedules under the EEA/EMEP
Guidebook (2023) framework (Section 1.A.3.a): fuel burn by engine type and
stage length times a fixed emission factor (3.15 kg CO2 per kg Jet A, 3.10
AvGas). Stage-resolved: airport x month x domestic/international, departure
(taxi-out, take-off, climb-out, cruise) and arrival (approach, taxi-in,
cruise) sides, plus flights, seats, seat-km. 1996-2024H1; partial July 2024
dropped; annual analysis ends 2023. Schedule-based (no load factors) and
aircraft-only [FANGYU: check description].
\paragraph{Allocation rules.} Three country aggregations from the stages:
LTO-territorial, fuel-uplift (cruise to departure; headline), 50/50 split.
Cruise recorded once per side, used once.
\paragraph{Validation.} 2019 world total 0.919 Gt vs ICCT 0.92 Gt; LTO share
11.8 per cent; log-log correlation with OWID/ICCT country series 0.99 in all
overlapping years; level ratio 1.05-1.10 (full-capacity scheduling); GACI
airport-year match 100 per cent after crosswalk patches.
\paragraph{Connectivity and sample.} GACI panel (Cheung et al.\ 2020 method,
1996-2024), country aggregations cwm/max/sum; estimation sample 184
countries, 1996-2023, N = 4,634.

%=======================================================================
\section{Empirical strategy}
%=======================================================================
% ln CO2_ct = beta ln GACI_cwm,ct + gamma ln pop_ct + delta ln seaMA_ct
%             + alpha_c + tau_t + eps_ct   (2SLS, robust SE)
\paragraph{Primary instrument.} Feyrer-type air/sea market-access shifter:
feyrer_int = a_t x ln(air market access, 1996 geography), with the sea-distance
market-access channel controlled directly (ln_sea_ma). KP F = 153.
% Validity notes for the text: survives interacting the cycle with baseline
% population, GDP, and capacity (F 120 -> 106), the stress test that kills
% naive alternatives (Appendix A); Conley breakdown f* = 0.85.
\paragraph{Secondary instrument (composition lens).} Tourism-heritage shifter
(natural+mixed endowment x global tourism cycle). Weaker and fragile to
size-x-cycle controls (F 18 -> 0; disclosed in Appendix A), but its compliers
are hub-quality changes, which isolates the composition margin and covers the
post-2010 era (the Feyrer instrument identifies mainly pre-2010: F 127 vs 7).
\paragraph{Airport level.} Descriptive mechanism evidence only (no valid
airport-level instrument: a heritage-proximity pilot died under size-x-cycle
controls, F 59 -> 2.3; noted in Appendix A). Causal language stays at the
country level. [YIFU: accident-exposure IV would upgrade this block.]

%=======================================================================
\section{Results}
%=======================================================================
\subsection{Main estimates: scale beats efficiency}
\begin{table}[H]
\centering
\caption{2SLS (Feyrer IV): connectivity and aviation CO2.}
\begin{threeparttable}
\small
\begin{tabular}{lcccccc}
\toprule
 & Total & LTO & 50/50 & Intl & Seat-km & Intensity \\
\midrule
$\ln\mathrm{GACI}_{cwm}$ & 5.669\sym{***} & 5.921\sym{***} & 5.620\sym{***} & 5.690\sym{***} & 6.068\sym{***} & $-0.399$\sym{***} \\
 & (0.412) & (0.448) & (0.409) & (0.415) & (0.443) & (0.118) \\
\midrule
Wald $p(\beta=1)$ & 0.000 & 0.000 & 0.000 & 0.000 & 0.000 & 0.000 \\
First-stage KP $F$ & 153.6 & 153.6 & 153.6 & 152.7 & 153.6 & 153.6 \\
Observations & 4{,}634 & 4{,}634 & 4{,}634 & 4{,}633 & 4{,}634 & 4{,}634 \\
\bottomrule
\end{tabular}
\begin{tablenotes}\footnotesize
\item Country and year FE; controls ln population and ln sea market access;
robust SE; instrument = air/sea market-access shifter. Intensity =
$\ln(\mathrm{CO}_2/\mathrm{seat\mbox{-}km})$.
\sym{***}~$p<0.01$.
\end{tablenotes}
\end{threeparttable}
\end{table}
% Reading: H1 supported with beta far above one, stable across allocation
% rules. Traffic (6.07) outruns emissions (5.67): the technique margin
% (-0.40) offsets roughly 0.4 log points but cannot overturn scale.
\subsection{Measures and timing}
% H2 aggregations, intl outcome: cwm 5.690 (0.415), sum 2.291 (0.237),
% max 4.836 (0.343), all p<0.01, F 90-171.
% Timing: Feyrer identifies pre-2010 (6.26, F 127); tourism identifies
% post-2010 (F 27.7). Complementary eras, same sign.
\subsection{The composition lens (tourism IV)}
% Under the hub-quality instrument: total emissions n.s. (1.16, SE 0.86) but
% international emissions 3.16 (SE 0.77), intensity +1.23 (SE 0.46):
% hub-quality changes recompose traffic toward international long-haul
% rather than scaling it. Present as H3 evidence with the instrument's
% fragility disclosed (Appendix A); flat Feyrer quintiles (appendix figure)
% show the scale response is uniform across income/connectivity/remoteness.
\subsection{Airport-level technique evidence}
% (descriptive; airport FE + country x year FE, SE clustered by airport)
% ln intensity -0.494 (0.044), intl share +0.331 (0.019); CO2 and seat-km
% columns (+3.7, +4.2) partly mechanical (GACI embeds capacity).
% Efficiency curve figure: 284 -> 77 kg per 1,000 seat-km across GACI
% ventiles, intl share 1 -> 60 per cent: the two curves cross.
% Hub table: DXB 19.7 Mt, LHR 18.6, LAX 16.4 top attributed airports.
\subsection{Attribution: what connectivity growth built}
% attributed share = 1 - exp(-beta dln GACI_cwm), applied to 2023 levels
% (same formula as the trade paper, cited as methodology only):
% 356 Mt = 42.5% of 2023 aviation CO2 (intl: 212 Mt, 41.2%).
% Country ranking: CHN +87.7, USA +36.4, ARE +24.7, JPN +17.4, TUR +16.6,
% IND +15.9, KOR +14.1; DEU -15.8 (declining connectivity).
% Caveat: beta >> 1 makes country-level attributed shares extreme for fast
% growers (~96%); world total moderate (42.5%); partial equilibrium.
% [TODO: attribution map figure]
\subsection{SAF scenarios: bending the curve}
% Fuel backed out as CO2/3.15; E(s) = E[(1-s) + s(1-r)], ReFuelEU path
% (2/6/20/34/42/70%, 2025-2050), life-cycle saving r in {50, 65, 80}%.
% On the 2023 total of 838 Mt: 2035 mandate -> 754/729/704 Mt; 2050 ->
% 545/457/369 Mt. Message: the technique margin improves with connectivity
% but cannot overturn scale (H1 vs H2); fuel switching is the lever that
% bends the level. [TODO: rebuild CO2_saf_scenarios.png in Mt units]

%=======================================================================
\section{The geography of attribution}
%=======================================================================
% Mismatch descriptives (CO2_map_levels_2023.png, CO2_map_mismatch_2023.png):
% 1. Round-trip symmetry nets endpoint mismatch (|dep-arr cruise| <= 0.7 Mt);
%    the binding margin is cruise attribution vs physical LTO.
% 2. Attributed-minus-physical world share 2023: USA +1.4 pp, ARE +1.4,
%    GBR +1.0, QAT +0.7, SGP +0.7 vs CHN -3.5, JPN/IDN/IND -0.5.
%    Airport decomposition: positive at every international mega-hub
%    (LHR +1.17, DXB +1.11), negative at large domestic hubs (SHA, XIY,
%    SZX, HGH, KMG, CKG, WUH).
% 3. Overlay with attribution (5.5): where connectivity-driven emissions
%    grew (China, Gulf) is not where they are booked under fuel-uplift
%    accounting: CORSIA/bunker-convention incidence discussion.
% [FANGYU: base-year bilateral route shares -> spillover regressions.]

%=======================================================================
\section{Discussion and roadmap}
%=======================================================================
% - Scale beats efficiency: connectivity policy cannot rely on hub
%   efficiency to decarbonise growth; SAF and demand-side measures price
%   the margin.
% - Accounting: attribution rules reallocate responsibility across
%   countries by up to 3.5 pp of the world total; CORSIA design implication.
% - Limitations: schedule-based emissions (no load factor), pre-2010
%   identification for the primary IV, no airport-level causal estimate yet,
%   partial-equilibrium scenarios, non-CO2 effects out of scope.
% - Roadmap: [YIFU] accident-exposure IV (airline-year and hub-airport-year
%   panel, ASN/ICAO ADREP 1996-2024, predetermined route-share exposure);
%   [FANGYU] bilateral base-year shares for network spillovers.

%=======================================================================
\appendix
\section{Instrument validity}
%=======================================================================
% (a) Size-x-cycle stress test (the discipline device of this paper):
%     first-stage F when a_t x baseline {pop, GDP, capacity} are controlled:
%       Feyrer: 120.1 -> 115.5 -> 105.8 (survives)
%       tourism: 13.8 -> 4.8/4.1 -> 0.0-0.6 (dies; hence secondary/lens)
%       airport heritage-proximity pilot: 59.5 -> 2.3 (abandoned)
% (b) Conley plausibly exogenous UCI (gamma in [0, gmax], share of RF):
%       intl CO2: f* = 0.85, UCI at 30% = [3.25, 4.67+]
%       total CO2: f* = 0.85; intensity: f* = 0.42
% (c) Flat Feyrer quintile reduced forms (CO2_rf_quintile appendix figure:
%     uniform ~0.7): scale response homogeneous; heterogeneity claims are
%     reserved for the composition lens.
% (d) [TODO: overidentification combining feyrer_int and tourism_int with
%     the LATE-difference discussion; H3 lens fragility disclosure.]
\end{document}
